\chapter{Results and Discussion}
\label{ch:results}

% RULE: every number in this chapter comes from experiments/results/*.csv;
% every figure regenerates via `make figures`. Never report the random split
% alone. Always include falling recall.

\section{Dataset Summary}
% Table: subjects, sessions, environments, windows per class. TODO after
% collection.
[TODO]

\section{Model Comparison Across Splits}

\figureorplaceholder{split_comparison.png}{Macro-F1 of each model across the
three evaluation splits (mean $\pm$ std over seeds/folds).}{fig:splits}

% Table: accuracy / macro-F1 / falling recall per model per split.
% Discuss the random-vs-cross-session gap (temporal leakage) explicitly —
% this is a strength of the evaluation, not a weakness of the model.
[TODO]

\section{Fall Detection}

\figureorplaceholder{falling_recall.png}{Falling recall (safety-critical
class) across splits.}{fig:falling}

[TODO]

\section{Ablation Studies}

\figureorplaceholder{ablation_receivers.png}{Effect of receiver count on
macro-F1.}{fig:abl-receivers}

\figureorplaceholder{ablation_window.png}{Effect of observation window
length on macro-F1.}{fig:abl-window}

\figureorplaceholder{ablation_rate.png}{Effect of packet rate on
macro-F1.}{fig:abl-rate}

[TODO]

\section{Real-Time Performance}
% End-to-end latency, sustained throughput, detection delay after activity
% onset. TODO measure on hardware.
[TODO]

\section{Failure Analysis}
% Confusion pairs (sitting vs lying), cross-subject drop, packet-loss
% sensitivity. Confusion matrices are auto-saved per run under docs/figures/.
[TODO]
